
%----------------------------------------------------------------------------------------
%	                            PORTADA
%----------------------------------------------------------------------------------------
% \makeportada{}{Daniel Carbajales Soto}{ 1° CFGS Automatización y Robótica Industrial}{}

%----------------------------------------------------------------------------------------
%                               TABLE OF CONTENTS
%----------------------------------------------------------------------------------------
% \maketoc
%----------------------------------------------------------------------------------------
%                                    HEADER 
%----------------------------------------------------------------------------------------
% \header{}{Daniel Carbajales Soto}{1° CFGSARI}
%----------------------------------------------------------------------------------------
%	                            TABLES
%----------------------------------------------------------------------------------------
% this is just a demo in case i need to make a table
%
%
%
% \begin{tabularx}{\linewidth}{|C{2cm}|Y|Y|C{2cm}|}
% \hline
% \makecell{Header 1} & \makecell{Header 2 \\ with linebreak} & \makecell{Header 3} & \makecell{Header 4} \\
% \hline
% Row 1 & This is a long cell text that automatically wraps in the Y column & Another long cell & Short \\
% \hline
% Row 2 & \multirow{2}{*}{Multirow text} & More text & Short \\
% \cline{1-1} \cline{3-4}
% Row 3 & & Extra text & Short \\
% \hline
% \end{tabularx}


%----------------------------------------------------------------------------------------
%	                            DOCUMENT
%----------------------------------------------------------------------------------------

\clearpage
% \fchapter{3}{Ej 7}
\phantomsection
\addcontentsline{toc}{chapter}{Ej 7}
\fsection{\boldit{ANÁLISIS DE SECUENCIA DE ACCIONES EN ENERGÍAS RENOVABLES.}}Una vez analizado cómo se conectan las SP,  elabora
una sucesión de acciones que podrián tener lugar si se realizara una inversión significativa en energías renovables en un país 
con un alto indice de desempleo y pobreza.


\fsubsection{Conexión con Políticas Laborales y ODS}
Las energías renovables se alinean con los ODS: 
\begin{itemize}
	\item  \boldit{1:} Fin de la pobreza
	\item  \boldit{7:} Energía asequible y no contaminante
	\item  \boldit{8:} Trabajo decente y crecimiento económico
\end{itemize}
Ya que generan empleo verde (hasta 12 millones globales según IRENA) y reducen desigualdades. 
Políticas laborales inclusivas (formación y equidad de género) amplifican el impacto, conectando 
inversión con creación de valor local.

\fsubsection{Secuencia de Acciones}
La inversión provoca  consecuencias socioeconómicas positivas
\begin{enumerate}
\item \textbf{Planificación y Financiamiento Inicial}: Gobiernos y organismos (e.g., Banco Mundial) asignan fondos (triplicando inversiones según OMM) para proyectos solares/eólicos, priorizando zonas pobres con subastas competitivas.
\item \textbf{Creación de Empleo Directo}: Instalación genera puestos en construcción y mantenimiento (e.g., 37-39\% empleos indirectos en eólica/solar), capacitando mano de obra local vía programas de formación, reduciendo desempleo en 10-20\% en regiones afectadas.
\item \textbf{Acceso Energético y Reducción de Pobreza}: Energía asequible electrifica hogares rurales (acelerando acceso como en datos del Banco Mundial), impulsando microempresas y agricultura productiva, bajando pobreza energética en 15-30\%.
\item \textbf{Crecimiento Económico e Innovación}: Cadena de suministro local fomenta industrias (e.g., fabricación de paneles), atrayendo IED y creando empleos indirectos (transporte, servicios), con ROI en 5-7 años vía ahorros en importaciones fósiles.
\item \textbf{Impactos Ambientales y Resiliencia}: Reducción de emisiones fortalece adaptación climática, liberando recursos para salud/educación, y promueve equidad (e.g., mujeres en roles técnicos).
\item \textbf{Evaluación y Escalabilidad}: Monitoreo mide impactos (empleo, PIB), ajustando políticas para sostenibilidad a largo plazo, como transiciones justas.
\end{enumerate}
En resumen, esta secuencia transforma desafíos en oportunidades, pero requiere gobernanza inclusiva para equidad.





%%%%%%%%%%%%%%%%%%%%%%%%%%%%%%%%%%%%%%%%%%%%%%%%%%%%%%%%%%
%%%%%%%%%%%%%%%%%%%% bibliography %%%%%%%%%%%%%%%%%%%%%%%%
% \nocite{*}                                 %%%%%%%%%%%%%
% \printbibliography                         %%%%%%%%%%%%%
%%%%%%%%%%%%%%%%%%%%%%%%%%%%%%%%%%%%%%%%%%%%%%%%%%%%%%%%%%
%%%%%% Last page in case there is no bibliography %%%%%%%%
% \label{Last Page}                           %%%%%%%%%%%%%
%%%%%%%%%%%%%%%%%%%%%%%%%%%%%%%%%%%%%%%%%%%%%%%%%%%%%%%%%%
